\section{Layer L0: identification and property dictionaries}
\label{sec:layer0}

Nothing an agent does is safe until it knows \emph{which thing} it is talking
about, and nothing it learns is transferable until two vendors calling the same
property by different names can be recognised as agreeing. L0 supplies both:
globally unique identifiers, and dictionaries of standardised properties bound
to those identifiers.

\subsection{IRDI and the ISO identification stack}

The International Registration Data Identifier (IRDI, ISO/IEC~11179-6 and
ISO~29002-5 \cite{iso29002}) is the identifier currency of European industrial
semantics. An IRDI such as \texttt{0173-1\#02-AAO677\#002} names a registration
authority (\texttt{0173-1} = ECLASS), an item type (\texttt{02} = property), an
item code and a \emph{version}. The versioning is the point: a property definition can
evolve without the identifier silently changing meaning, which is precisely the
guarantee a long-lived agent memory requires.

\paragraph{Agent affordance.} An IRDI is the join key across the entire stack.
The same identifier appears as an AAS \texttt{semanticId}, as an \opcua{}
\texttt{DictionaryEntry} reference, in an AutomationML attribute and in an
ECLASS-classified catalogue. An agent that carries IRDIs instead of English
labels can transfer a learned fact (``\emph{this} property is the one that must
match the recipe'') from one vendor's machine to another's. Without them, every
plant is a fresh vocabulary-learning problem. In the AUF this is the only
technology family scoring $3$ on U1 across the board.

\paragraph{Limitation.} IRDIs are not dereferenceable in the web sense by
default. An agent cannot simply \texttt{GET} one and receive a definition; it
needs a dictionary service. This is a small friction with a large practical
effect on retrieval-augmented pipelines, and it is why the survey scores
L0 lower on U8 than L2.

\subsection{ECLASS and IEC~CDD}

ECLASS \cite{eclass} is a hierarchical classification of products and services
with, for each classification class, a set of standardised properties, values
and value lists; ECLASS~Advanced additionally supplies structured property
blocks and aspects. IEC~CDD, the Common Data Dictionary built on IEC~61360
\cite{iec61360}, does the same for electrotechnical components, and IEC~61987
\cite{iec61987} for process measurement and control devices, each entry carrying
a preferred name, a definition, a data type, a unit and often a value list, in
several languages.

\paragraph{Agent affordance.} These dictionaries are, in effect,
\emph{pre-written ontologies of the properties an agent will encounter}, with
the three attributes that matter most for constrained generation: the datatype,
the unit and the closed value list for enumerated properties. When an
AAS submodel element carries an IEC~61360 data specification, an agent knows not
just that \texttt{FillVolume} is a number but that it is a
\texttt{REAL\_MEASURE} in millilitres. Our reference plant carries these as
\texttt{semanticId}s on every signal precisely to make this affordance
measurable; in Experiment~E1 it is part of what lifts C2 grounding to
\GrndTopOneCTwo{}.

\paragraph{Limitation.} L0 dictionaries describe \emph{property types}, never
\emph{instances}, and never relations between instances. They cannot say that
this filler is downstream of that pasteuriser, nor that this setpoint is
currently locked. They also carry licensing friction: ECLASS content is
commercially licensed, which materially complicates embedding it in an
open agent stack. This obstacle is under-discussed in the literature.

\subsection{ETIM, UNSPSC, GS1}

ETIM (technical product classification, strongest in electrical and building
trades) is structurally similar to ECLASS and increasingly harmonised with it.
UNSPSC is a procurement taxonomy: useful for commercial reasoning, nearly
useless for engineering, because it classifies without describing properties.
GS1 identifiers (GTIN for products, GIAI for individual assets, SSCC for
logistic units) together with EPCIS event vocabulary are the strongest L0 story
for \emph{instance} identity and for supply-chain provenance, and they are the
natural bridge when an agent's reasoning crosses from the plant floor to the
supply chain.

\subsection{What L0 provides to an agent, and what it does not}

L0 lets an agent map a phrase to a versioned, global property identifier,
recognise that two vendors' spellings denote the same property, and obtain the
datatype, unit and value list needed for constrained generation. It cannot
distinguish two \emph{instances} of the same property type, supply current
values, relations or state, or say whether acting is permitted now.

In our ablation, L0 alone corresponds to no experimental condition, because a
dictionary without an information model is not a deployable context. Its effect
is visible indirectly: it is the reason condition~C2 can express
\texttt{semanticId} and thereby recognise that the Line~1 and Line~2 fill-volume
setpoints are instances of \emph{the same property}, which is what makes the
recipe-conformance tasks (REC-01--03) tractable at all.
